%%%%%%%%%%%%Outline%%%%%%%%%%%%%%%%%%

%%%
% message<xx>
% xx

%%%%%%%%%%%%Outline%%%%%%%%%%%%%%%%%%

\section{Related Work}
\label{sec:related-work}

\textbf{BPF bytecode optimization.}
K2~\cite{k2} applies stochastic superoptimization to BPF bytecode.
Merlin~\cite{merlin} lifts BPF to LLVM IR for optimization.
EPSO~\cite{epso} uses evolutionary program synthesis at the bytecode
level. All three operate \emph{before} program loading: they require
the original \texttt{.bpf.o} file, cannot optimize programs already
loaded by third-party tools, and do not support runtime profiling
feedback. Since the BPF ISA lacks opcodes for conditional moves,
rotate, or wide-load fusion, no bytecode rewrite can introduce these
native instructions. \textsc{BpfReJIT} operates \emph{post-load} and
transparently, and is complementary to bytecode optimizers.

\textbf{BPF verification and extensibility.}
ePass~\cite{epass} proposes in-kernel optimization passes at BPF load
time but requires source modifications and does not support post-load
re-optimization. BCF~\cite{bcf} offloads BPF verification to userspace
via cryptographic certificates, addressing verifier scalability rather
than JIT code quality. \textsc{BpfReJIT} is orthogonal: BCF offloads
verification; \textsc{BpfReJIT} offloads optimization.

\textbf{Superoptimization.}
Massalin's superoptimizer~\cite{superoptimizer} and
STOKE~\cite{stoke} find optimal native sequences.
Bansal and Aiken~\cite{credal} automate peephole rule generation.
These systems discover better instruction sequences;
\textsc{BpfReJIT} provides a framework to \emph{deploy} such
sequences to kernel-resident programs transparently and safely.

\textbf{Feedback-directed and adaptive JIT compilation.}
Profile-guided optimization~\cite{feedback_directed} and adaptive JIT
techniques~\cite{jit_survey} (e.g., JVM HotSpot~\cite{jvm_hotspot},
V8, Graal) use runtime feedback to guide multi-tier compilation. These
systems have full control over their compiler backend.
\textsc{BpfReJIT} brings the same principle---runtime-guided dynamic
recompilation---to a much more constrained setting: the kernel BPF JIT,
where the compiler is part of the trusted computing base and cannot be
replaced, but can be extended through a minimal, safe interface.

\textbf{Auto-tuning and algorithm/schedule separation.}
Halide~\cite{halide} separates algorithm from schedule; TVM and XLA
auto-tune compilation parameters for deep learning workloads.
\textsc{BpfReJIT}'s microkernel-inspired separation of ``what can be
optimized'' from ``how to optimize'' draws on the same principle of
externalizing optimization decisions, applied to OS kernel compilation
with much stronger safety constraints.

\textbf{Compiler extensibility.}
LLVM's pass infrastructure~\cite{llvm} allows new optimization passes
to be registered as plugins. \textsc{BpfReJIT}'s kinsn mechanism
provides an analogous extensibility model for the kernel JIT: new
platform-specific instructions can be added via kernel modules without
modifying core kernel source, similar to how LLVM passes extend the
compiler without modifying the pass manager.

\textbf{Kernel live patching.}
Linux live patching (kpatch, livepatch) modifies running kernel code
at function granularity but must fully trust the patch author---there
is no verification step. \textsc{BpfReJIT} operates at BPF bytecode
granularity with full verifier re-validation, providing a fundamentally
stronger safety guarantee: the kernel verifier ensures that no
recompiled program can violate kernel safety invariants, regardless of
daemon behavior.